\section{Phần 2: Ứng Dụng Data Fitting vào Dữ Liệu Thực Tế}

\subsection{Giới Thiệu}

\textbf{Mục tiêu.} Phần này minh họa toàn bộ quy trình áp dụng các phương pháp data fitting đã
cài đặt ở Phần 1 vào một bộ dữ liệu thực tế. Nhóm không chỉ huấn luyện mô hình hồi quy mà còn xây
dựng pipeline tiền xử lý có thể tái lập, tránh rò rỉ dữ liệu, xử lý missing values, ngoại lai và
đa cộng tuyến trước khi đánh giá mô hình trên tập kiểm tra độc lập.

\textbf{Bài toán.} Với các thông tin quan sát về quỹ đạo, transit, khối lượng hành tinh và đặc
trưng của sao chủ, liệu có thể xây dựng mô hình toán học để dự đoán bán kính hành tinh ngoài Hệ
Mặt Trời hay không? Biến mục tiêu được chọn là:
\begin{equation}
    y = \texttt{pl\_rade},
\end{equation}
trong đó \texttt{pl\_rade} là bán kính hành tinh theo đơn vị bán kính Trái Đất.

\textbf{Quy ước.} Mức ý nghĩa dùng cho các kiểm định và chọn biến thống kê là
$\alpha = 0.05$. Các chỉ số RMSE, MAE và $R^2$ của mô hình tuyến tính chính được tính trên thang
$\log(1+\texttt{pl\_rade})$, vì target đã được log-transform trong pipeline. Khi cần diễn giải
trực tiếp theo bán kính hành tinh, báo cáo có thêm bảng dự báo đã biến đổi ngược về thang gốc.

Toàn bộ kết quả chính thức của Phần 2 được lưu trong thư mục \texttt{part2/real/}. Các file quan
trọng gồm:
\begin{itemize}
    \item \texttt{preprocessing\_metadata.json}: thông tin log-transform, Winsorization, MICE,
          VIF và các feature cuối cùng.
    \item \texttt{model\_results.json}, \texttt{model\_summary.csv}: hệ số, metric train/test và
          kết quả cross-validation của OLS, Ridge, Lasso.
    \item \texttt{advanced\_results.json}, \texttt{kernel\_ridge\_search.csv}: kết quả Kernel
          Ridge Regression với RBF kernel.
    \item Các hình EDA và diagnostic trong \texttt{part2/real/eda/} và
          \texttt{part2/real/diagnostics/}.
\end{itemize}

% ------------------------------------------------------------
\subsection{Hiểu Dữ Liệu (Data Understanding)}
% ------------------------------------------------------------

\subsubsection{Nguồn dữ liệu và mô tả bài toán}

Bộ dữ liệu được sử dụng là dữ liệu hành tinh ngoài Hệ Mặt Trời từ NASA Exoplanet Archive. File gốc
\texttt{raw.csv} sau khi tải về có khoảng 320 cột, bao gồm nhiều nhóm metadata, sai số đo lường,
cờ giới hạn và các đại lượng vật lý. Từ file gốc, nhóm tạo ra file gọn hơn:
\begin{center}
    \texttt{part2/data/planet.csv}.
\end{center}

Điểm quan trọng trong flow là quá trình tạo \texttt{planet.csv} và quá trình tiền xử lý trong
\texttt{DataPipeline} là hai bước khác nhau. Hai giai đoạn đầu tiên được thực hiện trước pipeline
để biến \texttt{raw.csv} thành \texttt{planet.csv}; sau đó pipeline mới đọc \texttt{planet.csv},
áp dụng domain restriction, chia train/test, xử lý missing/outlier và lọc VIF.

File \texttt{planet.csv} có 4636 dòng và 16 biến ban đầu. Sau khi áp dụng ràng buộc miền để tập
trung vào nhóm hành tinh đá/Super-Earth:
\begin{equation}
    \texttt{pl\_rade} \le 1.6,\qquad \texttt{pl\_bmasse} \le 10,
\end{equation}
số mẫu còn lại dùng cho EDA và mô hình hóa là 1084 dòng.

Lý do của bước lọc này đến từ cơ chế vật lý của dữ liệu. Động lực vật lý
(\emph{physical mechanics}) cấu trúc nên một hành tinh khí khổng lồ
(\emph{Gas Giant}) khác đáng kể so với một hành tinh đá
(\emph{Terrestrial}) hoặc Super-Earth. Nếu giữ nguyên toàn bộ tập dữ liệu, mặt phẳng hồi quy sẽ
phải đồng thời mô tả nhiều chế độ vật lý khác nhau, dễ bị nhiễu bởi hiện tượng phương sai không
đồng nhất (\emph{heteroskedasticity}) do chênh lệch quy mô quá lớn giữa các nhóm hành tinh.
Vì vậy, nhóm tối ưu hóa khả năng nội suy cục bộ của mô hình bằng cách chỉ giữ nhóm hành tinh đá và
Super-Earth theo ngưỡng phân rã Fulton (\emph{Fulton Gap}): bán kính không vượt quá
$1.6\,R_\oplus$ và khối lượng không vượt quá $10\,M_\oplus$. Thao tác này được thực hiện ở tầng dữ
liệu thô của pipeline, trước MICE và trước chuẩn hóa, để bộ nội suy MICE cũng như các tham số
mean/std chỉ học đặc trưng phân phối của riêng nhóm hành tinh này.

\begin{table}[H]
\centering
\caption{Quy mô dữ liệu qua các bước chính}
\label{tab:part2_data_size}
\begin{tabular}{@{}lrrp{0.38\textwidth}@{}}
\hline
\textbf{Giai đoạn} & \textbf{Số dòng} & \textbf{Số cột} & \textbf{Ghi chú} \\
\hline
\texttt{raw.csv} & 6610 & 320 & File tải từ NASA Exoplanet Archive, còn metadata và dòng mô tả. \\
\texttt{planet.csv} & 4636 & 16 & Tập biến đã được chọn lọc ban đầu. \\
Sau domain restriction & 1084 & 16 & Giữ các mẫu có \texttt{pl\_rade}$\le1.6$ và \texttt{pl\_bmasse}$\le10$. \\
Train sau tiền xử lý & 867 & 11 & Dùng để fit preprocessing và mô hình. \\
Test sau tiền xử lý & 217 & 11 & Chỉ dùng để đánh giá cuối cùng. \\
\hline
\end{tabular}
\end{table}

\subsubsection{Các biến trong dữ liệu}

\begin{table}[H]
\centering
\caption{Nhóm biến ban đầu trong \texttt{planet.csv}}
\label{tab:part2_variables}
\begin{tabular}{@{}p{0.22\textwidth}p{0.35\textwidth}p{0.34\textwidth}@{}}
\hline
\textbf{Nhóm biến} & \textbf{Tên biến} & \textbf{Ý nghĩa} \\
\hline
Quỹ đạo & \texttt{pl\_orbper}, \texttt{pl\_orbsmax}, \texttt{pl\_orbeccen}
& Chu kỳ quỹ đạo, bán trục lớn, độ lệch tâm. \\
Transit & \texttt{pl\_trandur}, \texttt{pl\_trandep}, \texttt{pl\_imppar}
& Thời lượng transit, độ sâu transit, impact parameter. \\
Năng lượng & \texttt{pl\_eqt}, \texttt{pl\_insol}
& Nhiệt độ cân bằng và thông lượng bức xạ nhận từ sao chủ. \\
Hành tinh & \texttt{pl\_bmasse}, \texttt{pl\_rade}
& Khối lượng và bán kính hành tinh. \\
Sao chủ & \texttt{st\_teff}, \texttt{st\_rad}, \texttt{st\_mass}, \texttt{st\_met}, \texttt{st\_logg}
& Nhiệt độ hiệu dụng, bán kính, khối lượng, kim loại tính và log-g của sao chủ. \\
Hệ sao & \texttt{sy\_dist}
& Khoảng cách từ Trái Đất đến hệ sao. \\
\hline
\end{tabular}
\end{table}

\begin{table}[H]
\centering
\caption{Một vài dòng dữ liệu mẫu từ \texttt{planet.csv} trước domain restriction}
\label{tab:part2_sample_data}
\resizebox{\textwidth}{!}{%
\begin{tabular}{@{}rrrrrrrr@{}}
\hline
\textbf{\texttt{pl\_orbper}} & \textbf{\texttt{pl\_orbsmax}} & \textbf{\texttt{pl\_trandur}} &
\textbf{\texttt{pl\_trandep}} & \textbf{\texttt{pl\_bmasse}} & \textbf{\texttt{st\_teff}} &
\textbf{\texttt{sy\_dist}} & \textbf{\texttt{pl\_rade}} \\
\hline
1.0039 & 0.0175 & 1.4562 & 0.04911 & 3.570 & 4971 & 820.905 & 1.710 \\
8.1724 & 0.0779 & 3.9020 & 0.08643 & 11.000 & 5705 & 1061.770 & 3.323 \\
6.2839 & 0.0687 & 4.1240 & 0.00357 & 0.437 & 6022 & 493.175 & 0.800 \\
3.1739 & 0.0464 & 3.5066 & 0.03030 & 10.100 & 6747 & 1318.050 & 3.150 \\
56.3585 & 0.2698 & 6.9270 & 0.25596 & 18.700 & 5446 & 962.888 & 4.541 \\
\hline
\end{tabular}}
\end{table}

Bảng mẫu cho thấy lý do cần ràng buộc miền trước khi mô hình hóa: một số hành tinh có bán kính và
khối lượng vượt xa nhóm hành tinh đá/Super-Earth. Nếu trộn chung các nhóm vật lý khác nhau, mô
hình tuyến tính dễ học một quan hệ trung bình kém ổn định.

\subsubsection{Khám phá dữ liệu ban đầu}

Các thống kê mô tả trên 1084 mẫu sau domain restriction được tóm tắt trong
Bảng~\ref{tab:part2_describe}. Dữ liệu có thang đo rất khác nhau: \texttt{pl\_orbper} dao động từ
0.1769 đến 199.6688 ngày, \texttt{pl\_insol} từ 0.144 đến 8385.9, còn \texttt{sy\_dist} từ
6.8693 đến 2712.58. Đây là tín hiệu rõ rằng cần biến đổi thang đo trước khi hồi quy.

\begin{table}[H]
\centering
\caption{Thống kê mô tả của một số biến số sau domain restriction}
\label{tab:part2_describe}
\resizebox{\textwidth}{!}{%
\begin{tabular}{@{}lrrrrrrr@{}}
\hline
\textbf{Biến} & \textbf{Count} & \textbf{Mean} & \textbf{Std} & \textbf{Min} & \textbf{Median} & \textbf{75\%} & \textbf{Max} \\
\hline
\texttt{pl\_orbper} & 1084 & 8.4750 & 13.7854 & 0.1769 & 4.7227 & 9.4638 & 199.6688 \\
\texttt{pl\_orbsmax} & 998 & 0.0679 & 0.0571 & 0.0050 & 0.0526 & 0.0850 & 0.6193 \\
\texttt{pl\_trandur} & 1060 & 2.7193 & 1.6231 & 0.2920 & 2.3765 & 3.3293 & 15.7220 \\
\texttt{pl\_trandep} & 1033 & 0.0405 & 0.0865 & 0.0012 & 0.0205 & 0.0354 & 1.6536 \\
\texttt{pl\_insol} & 981 & 520.5439 & 931.4663 & 0.1440 & 182.3750 & 518.4370 & 8385.9000 \\
\texttt{pl\_bmasse} & 1084 & 2.1245 & 1.2501 & 0.0364 & 2.1800 & 2.7600 & 10.0000 \\
\texttt{st\_teff} & 1084 & 5132.7154 & 897.5846 & 2566.0000 & 5382.0000 & 5815.0000 & 6681.0000 \\
\texttt{sy\_dist} & 1078 & 517.5306 & 400.8614 & 6.8693 & 446.8675 & 768.5703 & 2712.5800 \\
\texttt{pl\_rade} & 1084 & 1.2264 & 0.2577 & 0.3098 & 1.2600 & 1.4400 & 1.6000 \\
\hline
\end{tabular}}
\end{table}

\begin{figure}[H]
    \centering
    \includegraphics[width=\textwidth]{assets/part2/histograms.png}
    \caption{Histogram các biến số trước tiền xử lý. Nhiều biến có đuôi phải dài, đặc biệt là
    \texttt{pl\_orbper}, \texttt{pl\_insol}, \texttt{pl\_trandep} và \texttt{sy\_dist}.}
    \label{fig:part2_histograms}
\end{figure}

\begin{figure}[H]
    \centering
    \includegraphics[width=0.95\textwidth]{assets/part2/correlation_heatmap.png}
    \caption{Ma trận tương quan Pearson. Các nhóm biến quỹ đạo/năng lượng như
    \texttt{pl\_orbper}, \texttt{pl\_orbsmax}, \texttt{pl\_eqt}, \texttt{pl\_insol} có tương quan
    đáng chú ý, nên cần kiểm tra đa cộng tuyến bằng VIF.}
    \label{fig:part2_corr}
\end{figure}

\begin{figure}[H]
    \centering
    \includegraphics[width=\textwidth]{assets/part2/scatter_top5_target.png}
    \caption{Scatter plot giữa target \texttt{pl\_rade} và các biến có tương quan cao nhất.
    Quan hệ giữa target và predictors có xu hướng tuyến tính cục bộ nhưng vẫn còn nhiễu và ngoại lai.}
    \label{fig:part2_scatter}
\end{figure}

\subsubsection{Kiểm tra chất lượng dữ liệu}

Tương tự ví dụ tham khảo, nhóm kiểm tra bốn nhóm vấn đề: dữ liệu trùng lặp, dữ liệu khuyết, dữ liệu
nhiễu/ngoại lai và dữ liệu mâu thuẫn.

\paragraph{Dữ liệu trùng lặp.}
Trên 1084 mẫu sau domain restriction, số dòng trùng lặp bằng 0.

\begin{table}[H]
\centering
\caption{Kiểm tra duplicate sau domain restriction}
\label{tab:part2_duplicates}
\begin{tabular}{@{}lr@{}}
\hline
\textbf{Chỉ tiêu} & \textbf{Giá trị} \\
\hline
Số dòng kiểm tra & 1084 \\
Số dòng duplicate & 0 \\
\hline
\end{tabular}
\end{table}

\paragraph{Dữ liệu khuyết.}
Missing values xuất hiện ở nhiều nhóm biến khác nhau. Nếu xóa toàn bộ dòng chứa missing, dữ liệu sẽ
bị giảm đáng kể và dễ lệch mẫu; vì vậy pipeline dùng MICE ở bước tiền xử lý.

\begin{table}[H]
\centering
\caption{Các cột có missing values nhiều nhất trước tiền xử lý}
\label{tab:part2_missing_eda}
\begin{tabular}{@{}lrrl@{}}
\hline
\textbf{Cột} & \textbf{Missing count} & \textbf{Missing (\%)} & \textbf{Kiểu dữ liệu} \\
\hline
\texttt{pl\_orbeccen} & 145 & 13.376 & float64 \\
\texttt{pl\_insol} & 103 & 9.502 & float64 \\
\texttt{pl\_orbsmax} & 86 & 7.934 & float64 \\
\texttt{pl\_eqt} & 79 & 7.288 & float64 \\
\texttt{pl\_trandep} & 51 & 4.705 & float64 \\
\texttt{pl\_imppar} & 44 & 4.059 & float64 \\
\texttt{st\_met} & 35 & 3.229 & float64 \\
\texttt{pl\_trandur} & 24 & 2.214 & float64 \\
\texttt{sy\_dist} & 6 & 0.554 & float64 \\
\hline
\end{tabular}
\end{table}

\paragraph{Dữ liệu nhiễu và ngoại lai.}
Một số biến có khoảng giá trị rất rộng, chẳng hạn \texttt{pl\_insol} có max 8385.9 trong khi median
chỉ 182.375. \texttt{pl\_trandep} cũng có skewness ban đầu 10.1802. Những giá trị cực trị này có
thể là quan sát thật trong thiên văn học, nên nhóm không xóa dòng mà dùng log-transform và
Winsorization để giảm ảnh hưởng đòn bẩy.

\paragraph{Dữ liệu mâu thuẫn.}
Các biến chính như chu kỳ quỹ đạo, khối lượng, bán kính, nhiệt độ sao và khoảng cách đều nằm trong
miền vật lý dương sau khi chọn mẫu. Biến \texttt{pl\_imppar} có min $-0.13$ và max $1.2325$, có thể
đến từ sai số fitting trong dữ liệu quan trắc transit. Nhóm không xem đây là conflict data để xóa
cứng; thay vào đó bước Winsorization sẽ giới hạn ảnh hưởng của các giá trị cực trị trên tập train.

% ------------------------------------------------------------
\subsection{Chuẩn Bị Dữ Liệu (Data Preparation)}
% ------------------------------------------------------------

\subsubsection{Lựa chọn biến (Select Data)}

Quy trình từ \texttt{raw.csv} sang \texttt{planet.csv} là bước chọn biến thủ công trước khi chạy
\texttt{DataPipeline}. Mục tiêu của bước này là loại các cột không mang thông tin dự báo độc lập,
giảm nhiễu cho ma trận thiết kế và tránh để mô hình học từ các biến rò rỉ target.

\paragraph{Giai đoạn 1: loại bỏ siêu dữ liệu và sai số đo lường.}
Bộ dữ liệu gốc từ NASA Exoplanet Archive có hơn 200 cột thông tin. Tuy nhiên, nhiều cột trong số
này là metadata quan sát hoặc thông tin phụ của phép đo, không phải đặc trưng vật lý trực tiếp của
hành tinh. Ở bước lọc thô đầu tiên, nhóm loại các nhóm cột sau:
\begin{itemize}
    \item \textbf{Cột sai số đo lường}: các cột có hậu tố \texttt{err1}, \texttt{err2}, ví dụ
          \texttt{pl\_orbpererr1}, \texttt{pl\_radeerr2}, \texttt{st\_tefferr1}. Đây là độ bất
          định trên/dưới của phép đo, không phải giá trị vật lý trung tâm. Với bài toán hồi quy
          điểm, nhóm dùng giá trị ước lượng chính thay vì đưa cả khoảng sai số vào ma trận thiết kế.
    \item \textbf{Cột giới hạn}: các cột có hậu tố \texttt{lim}, ví dụ \texttt{pl\_orbperlim},
          \texttt{pl\_radelim}, \texttt{st\_masslim}. Các cờ này cho biết phép đo là giới hạn
          trên/dưới, dễ làm mô hình tuyến tính hiểu nhầm nếu dùng như biến số thông thường.
    \item \textbf{Cột tham chiếu và metadata phát hiện}: tên hành tinh/sao chủ, mã định danh
          \texttt{HD}, \texttt{HIP}, \texttt{TIC}, \texttt{Gaia}, phương pháp/năm/cơ sở phát hiện,
          link bài báo, kính thiên văn và thiết bị đo. Đây là thông tin mô tả quá trình con người
          quan sát dữ liệu, không trực tiếp quyết định bán kính hành tinh.
\end{itemize}

\paragraph{Giai đoạn 2: tinh chỉnh bằng domain knowledge và thống kê.}
Từ tập biến sau lọc thô, nhóm tiếp tục loại các cột có nguy cơ làm sai lệch hồi quy đa biến. Đây là
bước dùng kiến thức vật lý thiên văn để quyết định biến nào thật sự nên đi vào file
\texttt{planet.csv}.

\begin{table}[H]
\centering
\caption{Các nhóm biến bị loại trước khi tạo \texttt{planet.csv}}
\label{tab:part2_raw_to_planet_filter}
\resizebox{\textwidth}{!}{%
\begin{tabular}{@{}p{0.23\textwidth}p{0.24\textwidth}p{0.45\textwidth}@{}}
\hline
\textbf{Nhóm bị loại} & \textbf{Ví dụ} & \textbf{Lý do loại bỏ} \\
\hline
Rò rỉ dữ liệu khái niệm
& \texttt{pl\_radj}, \texttt{pl\_dens}, \texttt{pl\_ratror}
& \texttt{pl\_radj} là cùng bán kính hành tinh nhưng đổi đơn vị từ Trái Đất sang Sao Mộc.
\texttt{pl\_dens} chứa trực tiếp $R$ trong công thức
$\rho = M / (\frac{4}{3}\pi R^3)$, còn \texttt{pl\_ratror} là tỷ số $R_p/R_s$.
Giữ các biến này sẽ khiến mô hình suy ngược target thay vì học quan hệ dự báo độc lập. \\
\hline
Đa cộng tuyến gần hoàn hảo
& \texttt{pl\_bmassj}, \texttt{sy\_plx}
& \texttt{pl\_bmassj} chỉ là \texttt{pl\_bmasse} theo đơn vị Sao Mộc.
\texttt{sy\_plx} gần như là nghịch đảo của khoảng cách ($d \approx 1/p$).
Giữ đồng thời các cặp này làm ma trận thiết kế dễ suy biến và hệ số OLS kém ổn định. \\
\hline
Đại lượng dẫn xuất
& \texttt{st\_lum}, \texttt{st\_dens}
& Các biến này bị chi phối bởi nhiệt độ, khối lượng và bán kính sao:
$L \propto R^2T^4$ và $\rho_\star \propto M_\star/R_\star^3$. Ưu tiên biến gốc giúp mô hình dễ
diễn giải hơn và giảm trùng lặp thông tin. Riêng \texttt{st\_logg} được giữ tạm trong
\texttt{planet.csv} để pipeline đánh giá cùng các biến sao chủ, sau đó bị loại ở bước đầu pipeline. \\
\hline
Băng tần quang học
& \texttt{sy\_bmag}, \texttt{sy\_vmag}, \texttt{sy\_jmag}, \ldots
& Các độ sáng biểu kiến ở nhiều dải lọc thường tương quan mạnh với nhau và với khoảng cách quan
sát. Chúng dễ làm tăng VIF cục bộ nhưng không trực tiếp mô tả cấu trúc vật lý của hành tinh. \\
\hline
Biến hệ quy chiếu quan sát
& \texttt{pl\_tranmid}, \texttt{pl\_orbtper}, \texttt{sy\_pmra}, \texttt{sy\_pmdec}
& Bán kính hành tinh phụ thuộc vào vật chất và cấu trúc hệ sao, không phụ thuộc trực tiếp vào thời
điểm con người quan sát transit hay chuyển động riêng của hệ sao trên nền trời Trái Đất. \\
\hline
Biến thiếu quá nhiều
& \texttt{st\_age}, \texttt{st\_vsin}, \texttt{pl\_projobliq}, \texttt{pl\_trueobliq}, \texttt{pl\_occdep}
& Các thông số như tuổi sao, vận tốc quay sao hoặc độ nghiêng trục quay rất khó đo trong thực tế,
thường có tỷ lệ missing rất cao. Đưa vào MICE có thể khiến phần lớn giá trị là suy đoán và làm bóp
méo phương sai ban đầu. \\
\hline
\end{tabular}}
\end{table}

Để làm rõ hơn quyết định chọn/lọc cột, Bảng~\ref{tab:part2_column_decisions} trình bày các nhóm
cột cụ thể trong file gốc và lý do xử lý. Các cột có cùng hậu tố như \texttt{err1},
\texttt{err2}, \texttt{lim}, \texttt{reflink} được xử lý theo cùng một quy tắc vì chúng lặp lại
cho hầu hết các đại lượng vật lý.

{\scriptsize
\setlength{\tabcolsep}{3pt}
\begin{longtable}{@{}p{0.28\textwidth}p{0.12\textwidth}p{0.50\textwidth}@{}}
\caption{Quyết định xử lý chi tiết các nhóm cột từ \texttt{raw.csv}}
\label{tab:part2_column_decisions}\\
\hline
\textbf{Cột / pattern trong \texttt{raw.csv}} & \textbf{Quyết định} & \textbf{Giải thích} \\
\hline
\endfirsthead
\caption[]{Quyết định xử lý chi tiết các nhóm cột từ \texttt{raw.csv} (tiếp)}\\
\hline
\textbf{Cột / pattern trong \texttt{raw.csv}} & \textbf{Quyết định} & \textbf{Giải thích} \\
\hline
\endhead
\texttt{pl\_name}, \texttt{hostname}, \texttt{pl\_letter},
\texttt{hd\_name}, \texttt{hip\_name}, \texttt{tic\_id}, \texttt{gaia\_dr2\_id},
\texttt{gaia\_dr3\_id}
& Loại
& Đây là định danh hành tinh/sao chủ. Chúng không mang ý nghĩa vật lý liên tục để giải thích bán
kính. Nếu mã hóa trực tiếp, mô hình có thể học thuộc từng hệ sao thay vì học quy luật tổng quát. \\
\hline
\texttt{discoverymethod}, \texttt{disc\_year}, \texttt{disc\_refname},
\texttt{disc\_pubdate}, \texttt{disc\_locale}, \texttt{disc\_facility},
\texttt{disc\_telescope}, \texttt{disc\_instrument}
& Loại
& Đây là metadata của quá trình phát hiện. Chúng phản ánh lịch sử quan sát và năng lực thiết bị,
không phải cơ chế vật lý quyết định kích thước hành tinh. \\
\hline
\texttt{rv\_flag}, \texttt{pul\_flag}, \texttt{ptv\_flag}, \texttt{tran\_flag},
\texttt{ast\_flag}, \texttt{obm\_flag}, \texttt{micro\_flag}, \texttt{etv\_flag},
\texttt{ima\_flag}, \texttt{dkin\_flag}, \texttt{pl\_controv\_flag}, \texttt{ttv\_flag}
& Loại
& Các cờ này cho biết hành tinh được phát hiện bằng kỹ thuật nào hoặc có tranh cãi/biến thiên timing
hay không. Chúng dễ đưa thiên lệch quan sát vào mô hình, trong khi mục tiêu là dự đoán từ đại lượng
vật lý của hành tinh và sao chủ. \\
\hline
Mọi cột \texttt{*err1}, \texttt{*err2}
& Loại
& Đây là sai số trên/dưới của phép đo. Với hồi quy điểm, nhóm dùng giá trị trung tâm của đại lượng.
Đưa thêm sai số vào như feature có thể làm mô hình học chất lượng quan trắc thay vì học cấu trúc
vật lý. \\
\hline
Mọi cột \texttt{*lim}
& Loại
& Cờ giới hạn cho biết giá trị chỉ là upper/lower bound. Nếu trộn với biến vật lý liên tục, mô hình
tuyến tính dễ hiểu sai cờ đo lường này như một tín hiệu vật lý. \\
\hline
Mọi cột \texttt{*reflink}, \texttt{*\_systemref}
& Loại
& Đây là nguồn tham chiếu và hệ quy chiếu thời gian của phép đo. Chúng phục vụ truy xuất tài liệu,
không phục vụ nội suy bán kính hành tinh. \\
\hline
\texttt{sy\_snum}, \texttt{sy\_pnum}, \texttt{sy\_mnum}, \texttt{cb\_flag}
& Loại
& Số sao, số hành tinh, số mặt trăng hoặc cờ circumbinary là mô tả cấp hệ. Chúng có thể liên quan
đến bối cảnh hình thành, nhưng trong tập hiện tại tín hiệu chính cần học nằm ở quỹ đạo, transit,
khối lượng hành tinh và đặc trưng sao chủ. \\
\hline
\texttt{pl\_radj}, \texttt{pl\_dens}, \texttt{pl\_ratror}
& Loại
& Đây là nhóm rò rỉ khái niệm. \texttt{pl\_radj} là cùng target đổi đơn vị;
\texttt{pl\_dens} chứa bán kính trong công thức mật độ; \texttt{pl\_ratror} chứa trực tiếp tỷ số
bán kính hành tinh/sao. \\
\hline
\texttt{pl\_bmassj}, \texttt{sy\_plx}
& Loại
& \texttt{pl\_bmassj} trùng thông tin với \texttt{pl\_bmasse} nhưng khác đơn vị. \texttt{sy\_plx}
là thị sai, gần như nghịch đảo của \texttt{sy\_dist}. Giữ cả hai cặp sẽ làm tăng đa cộng tuyến. \\
\hline
\texttt{pl\_angsep}
& Loại
& Góc phân tách là đại lượng quan sát phụ thuộc vào khoảng cách hệ sao và hình học quan sát từ Trái
Đất. Nó không mô tả trực tiếp cấu trúc nội tại của hành tinh. \\
\hline
\texttt{pl\_orbincl}, \texttt{pl\_ratdor}, \texttt{pl\_occdep}, \texttt{pl\_rvamp}
& Loại
& Các biến này hoặc là hình học/quỹ đạo quan sát, hoặc phụ thuộc mạnh vào phương pháp đo và nhiều
trường hợp thiếu. Nhóm ưu tiên các biến transit ổn định hơn như \texttt{pl\_trandur},
\texttt{pl\_trandep}, \texttt{pl\_imppar}. \\
\hline
\texttt{pl\_tranmid}, \texttt{pl\_orbtper}, \texttt{pl\_orblper}
& Loại
& Đây là thời điểm hoặc góc pha quỹ đạo. Bán kính hành tinh không phụ thuộc trực tiếp vào thời điểm
con người ghi nhận transit hay epoch của cận điểm quỹ đạo. \\
\hline
\texttt{pl\_projobliq}, \texttt{pl\_trueobliq}, \texttt{st\_age}, \texttt{st\_vsin},
\texttt{st\_rotp}, \texttt{st\_radv}
& Loại
& Đây là các đại lượng khó đo, thường thiếu nhiều hoặc mô tả động học sao ở tầng phụ. Nếu đưa vào
MICE, phần lớn giá trị có nguy cơ là suy đoán, làm giảm độ tin cậy của phương sai và tương quan gốc. \\
\hline
\texttt{st\_spectype}, \texttt{st\_metratio}
& Loại
& \texttt{st\_spectype} là biến phân loại phức tạp cần mã hóa riêng; \texttt{st\_metratio} là thông
tin phụ cho cách biểu diễn kim loại tính. Nhóm giữ \texttt{st\_teff}, \texttt{st\_rad},
\texttt{st\_mass}, \texttt{st\_met}, \texttt{st\_logg} ở \texttt{planet.csv} để đại diện sao chủ. \\
\hline
\texttt{st\_lum}, \texttt{st\_dens}
& Loại
& Đây là đại lượng dẫn xuất từ các biến sao chủ gốc. Giữ chúng cùng \texttt{st\_teff},
\texttt{st\_rad}, \texttt{st\_mass} sẽ trùng thông tin và làm hệ số OLS khó diễn giải. \\
\hline
\texttt{rastr}, \texttt{ra}, \texttt{decstr}, \texttt{dec}, \texttt{glat}, \texttt{glon},
\texttt{elat}, \texttt{elon}, \texttt{sy\_pm}, \texttt{sy\_pmra}, \texttt{sy\_pmdec}
& Loại
& Đây là vị trí và chuyển động riêng trên bầu trời. Chúng mô tả hệ quy chiếu quan sát từ Trái Đất,
không quyết định trực tiếp bán kính hành tinh. \\
\hline
Các cột quang học \texttt{sy\_bmag}, \texttt{sy\_vmag}, \texttt{sy\_jmag},
\texttt{sy\_hmag}, \texttt{sy\_kmag}, \texttt{sy\_umag}, \texttt{sy\_gmag},
\texttt{sy\_rmag}, \texttt{sy\_imag}, \texttt{sy\_zmag}, \texttt{sy\_w1mag},
\texttt{sy\_w2mag}, \texttt{sy\_w3mag}, \texttt{sy\_w4mag}, \texttt{sy\_gaiamag},
\texttt{sy\_icmag}, \texttt{sy\_tmag}, \texttt{sy\_kepmag}
& Loại
& Các băng tần này đo độ sáng biểu kiến qua nhiều bộ lọc. Chúng tương quan mạnh với nhau và với
khoảng cách, dễ tạo nhiễu photometry và đa cộng tuyến cục bộ. \\
\hline
\texttt{pl\_nnotes}, \texttt{st\_nphot}, \texttt{st\_nrvc}, \texttt{st\_nspec},
\texttt{pl\_nespec}, \texttt{pl\_ntranspec}, \texttt{pl\_ndispec}
& Loại
& Đây là số lượng ghi chú hoặc số lượng chuỗi đo trong archive. Chúng phản ánh độ phong phú dữ liệu
quan sát, không phải cơ chế vật lý tạo nên bán kính hành tinh. \\
\hline
\end{longtable}
}

Ngược lại, \texttt{planet.csv} giữ lại 16 cột có ý nghĩa vật lý trực tiếp hơn. Bảng
\ref{tab:part2_kept_columns} tóm tắt vai trò của từng cột được giữ.

\begin{table}[H]
\centering
\caption{Các cột được giữ trong \texttt{planet.csv} và vai trò mô hình hóa}
\label{tab:part2_kept_columns}
\resizebox{\textwidth}{!}{%
\begin{tabular}{@{}p{0.20\textwidth}p{0.20\textwidth}p{0.52\textwidth}@{}}
\hline
\textbf{Cột giữ lại} & \textbf{Nhóm} & \textbf{Lý do giữ lại} \\
\hline
\texttt{pl\_orbper} & Quỹ đạo & Chu kỳ quỹ đạo phản ánh điều kiện động lực học quanh sao chủ. Được giữ ban đầu nhưng sau đó bị VIF loại khi thông tin trùng với \texttt{pl\_orbsmax}. \\
\texttt{pl\_orbsmax} & Quỹ đạo & Bán trục lớn mô tả khoảng cách quỹ đạo, ảnh hưởng đến bức xạ nhận được và điều kiện hình thành hành tinh. \\
\texttt{pl\_orbeccen} & Quỹ đạo & Độ lệch tâm mô tả hình dạng quỹ đạo, có thể liên quan đến lịch sử động lực của hệ. \\
\texttt{pl\_trandur} & Transit & Thời lượng transit chứa thông tin về hình học transit và kích thước/quỹ đạo tương đối. \\
\texttt{pl\_trandep} & Transit & Độ sâu transit liên quan đến tỷ lệ diện tích che khuất, nên có tín hiệu dự báo bán kính. \\
\texttt{pl\_imppar} & Transit & Impact parameter mô tả vị trí đường đi của hành tinh qua đĩa sao, giúp hiệu chỉnh tín hiệu transit. \\
\texttt{pl\_eqt} & Năng lượng & Nhiệt độ cân bằng đại diện môi trường nhiệt của hành tinh. Được giữ ban đầu nhưng sau đó bị VIF loại do trùng thông tin với quỹ đạo/nhiệt độ sao. \\
\texttt{pl\_insol} & Năng lượng & Thông lượng bức xạ nhận từ sao chủ, giữ lại vì phản ánh điều kiện năng lượng tác động lên hành tinh. \\
\texttt{pl\_bmasse} & Hành tinh & Khối lượng hành tinh là predictor vật lý quan trọng nhất để dự đoán bán kính. \\
\texttt{st\_teff} & Sao chủ & Nhiệt độ sao chủ ảnh hưởng đến bức xạ và môi trường quanh hành tinh. \\
\texttt{st\_rad} & Sao chủ & Bán kính sao chủ liên quan trực tiếp đến diễn giải tín hiệu transit. \\
\texttt{st\_mass} & Sao chủ & Khối lượng sao chủ được giữ tạm để pipeline đánh giá, sau đó bị loại ở bước đầu do chồng lấp thông tin. \\
\texttt{st\_met} & Sao chủ & Kim loại tính có thể liên quan đến thành phần vật chất và lịch sử hình thành hệ hành tinh. \\
\texttt{st\_logg} & Sao chủ & Trọng lực bề mặt sao được giữ tạm để pipeline đánh giá, sau đó bị loại ở bước đầu vì là đại lượng dẫn xuất/chồng lấp. \\
\texttt{sy\_dist} & Hệ sao & Khoảng cách tới hệ sao giúp mô tả thiên lệch quan sát và độ tin cậy tín hiệu. \\
\texttt{pl\_rade} & Target & Bán kính hành tinh theo đơn vị Trái Đất, là biến mục tiêu cần dự đoán. \\
\hline
\end{tabular}}
\end{table}

Sau hai giai đoạn này, \texttt{planet.csv} là tập dữ liệu đã được chọn lọc trước pipeline, gồm các
biến vật lý có ý nghĩa hơn cho bài toán dự đoán \texttt{pl\_rade}. Cụ thể, file này giữ các nhóm
biến quỹ đạo, transit, năng lượng, khối lượng hành tinh, đặc trưng sao chủ và khoảng cách hệ sao.

\paragraph{Phân biệt với bước VIF trong pipeline.}
VIF không thuộc quá trình tạo \texttt{planet.csv}. VIF chỉ được tính sau khi pipeline đã đọc
\texttt{planet.csv}, áp dụng domain restriction, chia train/test, log-transform, Winsorization và
MICE imputation. Nói cách khác, giai đoạn 1--2 là chọn biến trước pipeline; còn VIF là thanh lọc
tự động bên trong pipeline trên tập train đã xử lý.

Trong pipeline chính, hai biến \texttt{st\_mass} và \texttt{st\_logg} được loại ở bước đầu vì thông
tin chồng lấp mạnh với các biến sao chủ còn lại. Sau đó VIF tiếp tục loại các biến đa cộng tuyến
vượt ngưỡng như trình bày ở Bảng~\ref{tab:part2_vif_history}.

\subsubsection{Làm sạch dữ liệu (Clean Data)}

Pipeline tuân thủ nguyên tắc tránh data leakage: chia train/test trước, sau đó chỉ fit tham số tiền
xử lý trên train và áp dụng lại cho test.

\begin{align}
    &\text{Load data}
    \rightarrow
    \text{Domain restriction}
    \rightarrow
    \text{Train/test split}, \nonumber\\
    &\text{Fit preprocessing trên train}
    \rightarrow
    \text{Transform test bằng tham số train}.
\end{align}

\begin{lstlisting}[language=Python, caption={Tách train/test trước khi fit preprocessing}]
raw = raw[(raw["pl_rade"] <= 1.6) & (raw["pl_bmasse"] <= 10)]
train_raw, test_raw = train_test_split_frame(
    raw, test_size=0.2, random_state=RANDOM_STATE
)

pipeline = DataPipeline(target="pl_rade")
X_train, y_train = pipeline.fit_transform(train_raw)
X_test,  y_test  = pipeline.transform(test_raw, include_target=True)
\end{lstlisting}

Sau khi chia dữ liệu, train có 867 dòng và test có 217 dòng. Tất cả các bước log-transform,
Winsorization, MICE, VIF và standardization đều được fit trên train.

\subsubsection{Xây dựng dữ liệu (Construct Data)}

\paragraph{Log-transform.}
Nhiều đại lượng thiên văn có quan hệ lũy thừa, ví dụ Định luật Kepler hoặc quan hệ bức xạ. Với
quan hệ:
\begin{equation}
    y = Cx^\alpha,
\end{equation}
lấy log hai vế đưa bài toán về dạng gần tuyến tính:
\begin{equation}
    \log y = \log C + \alpha \log x.
\end{equation}

Pipeline dùng phép biến đổi:
\begin{equation}
    z = \log(1+x),
\end{equation}
và target được chuyển thành:
\begin{equation}
    y = \log(1+\texttt{pl\_rade}) = \texttt{log\_pl\_rade}.
\end{equation}

\begin{table}[H]
\centering
\caption{Ảnh hưởng của log-transform lên skewness}
\label{tab:part2_log_transform}
\resizebox{\textwidth}{!}{%
\begin{tabular}{@{}llrrrrr@{}}
\hline
\textbf{Biến gốc} & \textbf{Biến mới} & \textbf{Missing} & \textbf{Min trước} & \textbf{Max trước} &
\textbf{Skew trước} & \textbf{Skew sau} \\
\hline
\texttt{pl\_orbper} & \texttt{log\_pl\_orbper} & 0 & 0.1769 & 199.6688 & 6.8615 & 0.4901 \\
\texttt{pl\_orbsmax} & \texttt{log\_pl\_orbsmax} & 73 & 0.0050 & 0.6193 & 3.2313 & 2.5948 \\
\texttt{pl\_trandep} & \texttt{log\_pl\_trandep} & 40 & 0.0012 & 1.6536 & 10.1802 & 6.9522 \\
\texttt{pl\_insol} & \texttt{log\_pl\_insol} & 83 & 0.1440 & 8385.9000 & 3.9596 & -0.3815 \\
\texttt{pl\_bmasse} & \texttt{log\_pl\_bmasse} & 0 & 0.0364 & 10.0000 & 1.7694 & -0.2716 \\
\texttt{pl\_eqt} & \texttt{log\_pl\_eqt} & 67 & 171.7000 & 3186.0000 & 0.8788 & -0.4766 \\
\texttt{sy\_dist} & \texttt{log\_sy\_dist} & 6 & 6.8693 & 2712.5800 & 1.1829 & -1.1509 \\
\texttt{pl\_rade} & \texttt{log\_pl\_rade} & 0 & 0.3098 & 1.6000 & -0.6930 & -0.9918 \\
\hline
\end{tabular}}
\end{table}

\begin{figure}[H]
    \centering
    \includegraphics[width=\textwidth]{assets/part2/log_transform_before_after.png}
    \caption{So sánh phân phối trước và sau log-transform. Các biến có đuôi phải dài được nén
    thang đo rõ rệt, đặc biệt là \texttt{pl\_orbper}, \texttt{pl\_insol} và \texttt{pl\_bmasse}.}
    \label{fig:part2_log_transform}
\end{figure}

\paragraph{Winsorization.}
Nhóm không xóa ngoại lai mà cap giá trị ở hai đầu phân phối theo bách phân vị 1\% và 99\%:
\begin{equation}
    x' =
    \begin{cases}
        q_{0.01}, & x < q_{0.01},\\
        x,        & q_{0.01} \le x \le q_{0.99},\\
        q_{0.99}, & x > q_{0.99}.
    \end{cases}
\end{equation}

\begin{table}[H]
\centering
\caption{Tóm tắt Winsorization trên tập train}
\label{tab:part2_winsor}
\begin{tabular}{@{}lrrrrr@{}}
\hline
\textbf{Cột} & \textbf{p01} & \textbf{p99} & \textbf{Dưới p01} & \textbf{Trên p99} & \textbf{Clipped (\%)} \\
\hline
\texttt{pl\_orbeccen} & 0.0000 & 0.2994 & 0 & 8 & 1.061 \\
\texttt{pl\_trandur} & 0.6241 & 8.1895 & 9 & 9 & 2.120 \\
\texttt{pl\_imppar} & 0.0040 & 0.9395 & 9 & 9 & 2.179 \\
\texttt{st\_rad} & 0.1963 & 1.7887 & 9 & 9 & 2.076 \\
\hline
\end{tabular}
\end{table}

\begin{figure}[H]
    \centering
    \includegraphics[width=\textwidth]{assets/part2/winsorization_before_after.png}
    \caption{Boxplot trước và sau Winsorization. Các điểm cực trị được kéo về ngưỡng fit trên
    train, trong khi phần trung tâm của phân phối được giữ nguyên.}
    \label{fig:part2_winsor}
\end{figure}

\paragraph{MICE imputation.}
Sau log-transform và Winsorization, tập train vẫn còn missing values ở 9 cột. Nhóm dùng MICE
(Multiple Imputation by Chained Equations) thay cho median imputation đơn giản để bảo toàn cấu trúc
tương quan đa biến.

\begin{table}[H]
\centering
\caption{Missing values sau log-transform trên tập train}
\label{tab:part2_missing_after_transform}
\begin{tabular}{@{}lrrl@{}}
\hline
\textbf{Cột} & \textbf{Missing count} & \textbf{Missing (\%)} & \textbf{Kiểu dữ liệu} \\
\hline
\texttt{pl\_orbeccen} & 113 & 13.033 & float64 \\
\texttt{log\_pl\_insol} & 83 & 9.573 & float64 \\
\texttt{log\_pl\_orbsmax} & 73 & 8.420 & float64 \\
\texttt{log\_pl\_eqt} & 67 & 7.728 & float64 \\
\texttt{pl\_imppar} & 41 & 4.729 & float64 \\
\texttt{log\_pl\_trandep} & 40 & 4.614 & float64 \\
\texttt{st\_met} & 24 & 2.768 & float64 \\
\texttt{pl\_trandur} & 18 & 2.076 & float64 \\
\texttt{log\_sy\_dist} & 6 & 0.692 & float64 \\
\hline
\end{tabular}
\end{table}

MICE trong pipeline dùng 5 bản imputation, 5 vòng chained iterations, hồi quy Ridge phụ với
regularization 0.001 và noise scale 0.15.

\begin{table}[H]
\centering
\caption{Tóm tắt giá trị được MICE impute trên tập train}
\label{tab:part2_mice}
\resizebox{\textwidth}{!}{%
\begin{tabular}{@{}lrrrrrr@{}}
\hline
\textbf{Cột} & \textbf{Missing} & \textbf{Missing (\%)} & \textbf{Median khởi tạo} &
\textbf{Observed mean} & \textbf{Imputed mean} & \textbf{Imputed std} \\
\hline
\texttt{pl\_orbeccen} & 113 & 13.033 & 0.000000 & 0.013696 & 0.033514 & 0.021444 \\
\texttt{log\_pl\_insol} & 83 & 9.573 & 5.207544 & 5.036231 & 5.332386 & 1.555943 \\
\texttt{log\_pl\_orbsmax} & 73 & 8.420 & 0.052213 & 0.064713 & 0.045031 & 0.040655 \\
\texttt{log\_pl\_eqt} & 67 & 7.728 & 6.856987 & 6.806412 & 6.805654 & 0.457823 \\
\texttt{pl\_imppar} & 41 & 4.729 & 0.284700 & 0.346475 & 0.394104 & 0.059489 \\
\texttt{log\_pl\_trandep} & 40 & 4.614 & 0.020498 & 0.036569 & 0.046997 & 0.040054 \\
\texttt{st\_met} & 24 & 2.768 & 0.000000 & -0.025205 & -0.097143 & 0.039923 \\
\texttt{pl\_trandur} & 18 & 2.076 & 2.372000 & 2.699806 & 2.127761 & 1.064210 \\
\texttt{log\_sy\_dist} & 6 & 0.692 & 6.104394 & 5.830641 & 5.143891 & 0.991066 \\
\hline
\end{tabular}}
\end{table}

\begin{figure}[H]
    \centering
    \includegraphics[width=\textwidth]{assets/part2/mice_observed_vs_imputed.png}
    \caption{So sánh phân phối observed và phần được MICE impute. Giá trị impute không bị gom
    về một điểm trung vị duy nhất mà vẫn có độ phân tán theo quan hệ đa biến.}
    \label{fig:part2_mice}
\end{figure}

\paragraph{Giai đoạn 3 trong DataPipeline: loại biến bằng VIF.}
Khác với hai giai đoạn tạo \texttt{planet.csv} ở trên, bước VIF diễn ra bên trong
\texttt{DataPipeline}. Sau khi train/test đã được tách và tập train đã đi qua log-transform,
Winsorization và MICE, pipeline kiểm tra đa cộng tuyến bằng hệ số phóng đại phương sai:
\begin{equation}
    \mathrm{VIF}_j = \frac{1}{1-R_j^2},
\end{equation}
trong đó $R_j^2$ là hệ số xác định khi hồi quy biến $x_j$ theo các biến còn lại. Ngưỡng loại biến
là $\mathrm{VIF} > 10$.

\begin{table}[H]
\centering
\caption{Lịch sử lọc VIF}
\label{tab:part2_vif_history}
\begin{tabular}{@{}lrl@{}}
\hline
\textbf{Biến có VIF cao nhất} & \textbf{VIF} & \textbf{Quyết định} \\
\hline
\texttt{log\_pl\_orbper} & 43.1374 & Drop \\
\texttt{log\_pl\_eqt} & 20.4514 & Drop \\
\texttt{st\_teff} & 9.2570 & Dừng vì nhỏ hơn 10 \\
\hline
\end{tabular}
\end{table}

\begin{table}[H]
\centering
\caption{VIF cuối cùng của các feature còn lại}
\label{tab:part2_final_vif}
\begin{tabular}{@{}lrlr@{}}
\hline
\textbf{Feature} & \textbf{VIF} & \textbf{Feature} & \textbf{VIF} \\
\hline
\texttt{st\_teff} & 9.2570 & \texttt{log\_pl\_trandep} & 1.8730 \\
\texttt{st\_rad} & 5.9915 & \texttt{pl\_orbeccen} & 1.2736 \\
\texttt{log\_pl\_insol} & 5.8778 & \texttt{st\_met} & 1.1861 \\
\texttt{log\_pl\_orbsmax} & 4.6054 & \texttt{pl\_imppar} & 1.0732 \\
\texttt{pl\_trandur} & 3.0931 & \texttt{log\_pl\_bmasse} & 1.0598 \\
\texttt{log\_sy\_dist} & 2.8147 & & \\
\hline
\end{tabular}
\end{table}

\begin{figure}[H]
    \centering
    \includegraphics[width=\textwidth]{assets/part2/vif_before_after.png}
    \caption{VIF trước và sau lọc. Việc loại \texttt{log\_pl\_orbper} và \texttt{log\_pl\_eqt}
    đưa tất cả VIF còn lại xuống dưới ngưỡng 10.}
    \label{fig:part2_vif}
\end{figure}

Sau VIF filtering, 11 feature cuối cùng là:
\begin{center}
\begin{tabular}{@{}lll@{}}
\hline
\texttt{pl\_orbeccen} & \texttt{pl\_trandur} & \texttt{pl\_imppar} \\
\texttt{st\_teff} & \texttt{st\_rad} & \texttt{st\_met} \\
\texttt{log\_pl\_orbsmax} & \texttt{log\_pl\_trandep} & \texttt{log\_pl\_insol} \\
\texttt{log\_pl\_bmasse} & \texttt{log\_sy\_dist} & \\
\hline
\end{tabular}
\end{center}

\paragraph{Chuẩn hóa dữ liệu.}
Tất cả feature cuối cùng được chuẩn hóa bằng z-score:
\begin{equation}
    x_{ij}^{\mathrm{std}} =
    \frac{x_{ij} - \mu_j^{\mathrm{train}}}{\sigma_j^{\mathrm{train}}},
\end{equation}
trong đó $\mu_j^{\mathrm{train}}$ và $\sigma_j^{\mathrm{train}}$ chỉ được fit trên train.

% ------------------------------------------------------------
\subsection{Xây Dựng Mô Hình (Modeling)}
% ------------------------------------------------------------

Bốn mô hình tuyến tính chính được so sánh:
\begin{enumerate}
    \item \textbf{OLS full}: dùng toàn bộ 11 feature sau tiền xử lý.
    \item \textbf{OLS selected}: chọn feature có p-value nhỏ hơn 0.05 từ OLS full.
    \item \textbf{Ridge}: dùng hàm \texttt{ridge\_fit} ở Phần 1 và chọn $\lambda$ bằng 5-fold CV.
    \item \textbf{Lasso}: dùng hàm \texttt{lasso\_fit} ở Phần 1 và chọn $\lambda$ bằng 3-fold CV.
\end{enumerate}

\subsubsection{OLS full và chọn biến theo p-value}

OLS full có $R^2_{\mathrm{train}} = 0.681784$ và
$R^2_{\mathrm{adj,train}} = 0.677690$. Các biến có ý nghĩa thống kê ở mức 0.05 được giữ lại cho
OLS selected:
\begin{center}
\texttt{pl\_trandur}, \texttt{pl\_imppar}, \texttt{st\_teff},
\texttt{log\_pl\_trandep}, \texttt{log\_pl\_bmasse}, \texttt{log\_sy\_dist}.
\end{center}

\begin{table}[H]
\centering
\caption{Bảng suy diễn hệ số OLS full}
\label{tab:part2_ols_inference}
\resizebox{\textwidth}{!}{%
\begin{tabular}{@{}lrrrrrr@{}}
\hline
\textbf{Feature} & \textbf{Coef} & \textbf{Std.Err.} & \textbf{t-stat} & \textbf{p-value} &
\textbf{CI lower} & \textbf{CI upper} \\
\hline
\texttt{pl\_orbeccen} & -0.004508 & 0.002704 & -1.6673 & 0.095816 & -0.009815 & 0.000799 \\
\texttt{pl\_trandur} & -0.011414 & 0.004214 & -2.7088 & 0.006888 & -0.019684 & -0.003144 \\
\texttt{pl\_imppar} & 0.005751 & 0.002482 & 2.3169 & 0.020743 & 0.000879 & 0.010622 \\
\texttt{st\_teff} & -0.014611 & 0.007290 & -2.0043 & 0.045348 & -0.028918 & -0.000303 \\
\texttt{st\_rad} & 0.009279 & 0.005864 & 1.5822 & 0.113964 & -0.002231 & 0.020790 \\
\texttt{st\_met} & 0.000139 & 0.002609 & 0.0533 & 0.957505 & -0.004982 & 0.005260 \\
\texttt{log\_pl\_orbsmax} & 0.007766 & 0.005142 & 1.5104 & 0.131315 & -0.002326 & 0.017857 \\
\texttt{log\_pl\_trandep} & 0.013692 & 0.003279 & 4.1757 & $3.28\times10^{-5}$ & 0.007256 & 0.020128 \\
\texttt{log\_pl\_insol} & -0.005905 & 0.005809 & -1.0165 & 0.309668 & -0.017305 & 0.005496 \\
\texttt{log\_pl\_bmasse} & 0.095317 & 0.002466 & 38.6447 & $8.78\times10^{-190}$ & 0.090476 & 0.100158 \\
\texttt{log\_sy\_dist} & 0.034443 & 0.004020 & 8.5689 & $4.84\times10^{-17}$ & 0.026554 & 0.042332 \\
\hline
\end{tabular}}
\end{table}

Biến \texttt{log\_pl\_bmasse} có hệ số dương lớn nhất và p-value cực nhỏ, phù hợp với trực giác
vật lý: hành tinh có khối lượng lớn thường có bán kính lớn hơn.

\subsubsection{Chọn siêu tham số cho Ridge và Lasso}

Ridge quét lưới $\lambda$ từ $10^{-3}$ đến $10^3$ theo thang log. Giá trị tốt nhất theo CV-MSE là:
\begin{equation}
    \lambda^*_{\mathrm{Ridge}} = 10.
\end{equation}

\begin{table}[H]
\centering
\caption{Một số kết quả Ridge CV}
\label{tab:part2_ridge_cv}
\begin{tabular}{@{}rrrr@{}}
\hline
\textbf{$\lambda$} & \textbf{Mean CV-MSE} & \textbf{Std CV-MSE} & \textbf{Mean CV-$R^2$} \\
\hline
0.0010 & 0.005075 & 0.001660 & 0.673407 \\
0.1000 & 0.005075 & 0.001659 & 0.673410 \\
1.0000 & 0.005075 & 0.001656 & 0.673433 \\
3.1623 & 0.005074 & 0.001650 & 0.673469 \\
10.0000 & 0.005074 & 0.001629 & 0.673433 \\
31.6228 & 0.005090 & 0.001571 & 0.672223 \\
100.0000 & 0.005240 & 0.001432 & 0.661976 \\
1000.0000 & 0.008594 & 0.001069 & 0.441426 \\
\hline
\end{tabular}
\end{table}

Sau khi mở rộng lưới, Lasso quét 13 giá trị
$\lambda \in [10^{-4}, 10^2]$ theo thang log bằng 3-fold CV. Lưới mới cho thấy vùng
regularization tốt khá phẳng ở các $\lambda$ nhỏ, sau đó CV-MSE bắt đầu tăng rõ khi
$\lambda \ge 0.3162$. Giá trị tốt nhất là:
\begin{equation}
    \lambda^*_{\mathrm{Lasso}} = 0.0316228.
\end{equation}

\begin{table}[H]
\centering
\caption{Kết quả Lasso CV}
\label{tab:part2_lasso_cv}
\resizebox{\textwidth}{!}{%
\begin{tabular}{@{}rrrr@{}}
\hline
\textbf{$\lambda$} & \textbf{Mean CV-MSE} & \textbf{Std CV-MSE} & \textbf{Mean CV-$R^2$} \\
\hline
0.0001 & 0.005093 & 0.001421 & 0.672720 \\
0.0003 & 0.005093 & 0.001421 & 0.672720 \\
0.001 & 0.005093 & 0.001421 & 0.672720 \\
0.003 & 0.005093 & 0.001421 & 0.672722 \\
0.010 & 0.005093 & 0.001420 & 0.672726 \\
0.032 & 0.005092 & 0.001420 & 0.672730 \\
0.100 & 0.005093 & 0.001418 & 0.672675 \\
0.316 & 0.005106 & 0.001415 & 0.671822 \\
1.000 & 0.005140 & 0.001420 & 0.669538 \\
3.162 & 0.005301 & 0.001387 & 0.658832 \\
10.000 & 0.005953 & 0.001265 & 0.615589 \\
31.623 & 0.008702 & 0.000978 & 0.434213 \\
100.000 & 0.015532 & 0.001217 & -0.012201 \\
\hline
\end{tabular}}
\end{table}

Với $\lambda^* = 0.0316228$, Lasso vẫn không triệt tiêu hệ số nào: $0/11$ feature có hệ số bằng
0 theo ngưỡng $10^{-10}$. Điều này cho thấy Lasso trong lần chạy này hoạt động giống một OLS được
làm mượt nhẹ hơn là một bộ chọn biến thưa. Lưới mở rộng cũng làm rõ rằng nếu phạt mạnh hơn
($\lambda \ge 1$), sai số CV tăng dần; do đó việc không chọn nghiệm thưa là kết quả thực nghiệm
của dữ liệu, không phải do lưới ban đầu quá hẹp.

% ------------------------------------------------------------
\subsection{Dự Báo (Prediction)}
% ------------------------------------------------------------

Mô hình tuyến tính tốt nhất trên test set là Lasso. Vì target được huấn luyện trên thang
\texttt{log\_pl\_rade}, dự báo được biến đổi ngược bằng:
\begin{equation}
    \widehat{\texttt{pl\_rade}} = \exp(\hat{y}) - 1.
\end{equation}

\begin{table}[H]
\centering
\caption{Một số dự báo của mô hình Lasso trên tập test}
\label{tab:part2_prediction_sample}
\resizebox{\textwidth}{!}{%
\begin{tabular}{@{}rrrrrrr@{}}
\hline
\textbf{STT} & \textbf{Actual log} & \textbf{Pred log} & \textbf{Residual log} &
\textbf{Actual radius} & \textbf{Pred radius} & \textbf{APE (\%)} \\
\hline
1 & 0.703098 & 0.673404 & 0.029694 & 1.020000 & 0.960901 & 5.794 \\
2 & 0.774727 & 0.788004 & -0.013276 & 1.170000 & 1.199002 & 2.479 \\
3 & 0.904218 & 0.850695 & 0.053524 & 1.470000 & 1.341272 & 8.757 \\
4 & 0.751416 & 0.762798 & -0.011382 & 1.120000 & 1.144267 & 2.167 \\
5 & 0.797507 & 0.797699 & -0.000192 & 1.220000 & 1.220426 & 0.035 \\
6 & 0.548121 & 0.601355 & -0.053234 & 0.730000 & 0.824590 & 12.958 \\
7 & 0.581098 & 0.616384 & -0.035286 & 0.788000 & 0.852218 & 8.150 \\
8 & 0.943906 & 0.879492 & 0.064414 & 1.570000 & 1.409674 & 10.212 \\
9 & 0.783902 & 0.787921 & -0.004019 & 1.190000 & 1.198819 & 0.741 \\
10 & 0.662688 & 0.647882 & 0.014805 & 0.940000 & 0.911489 & 3.033 \\
\hline
\end{tabular}}
\end{table}

Trên thang bán kính gốc, Lasso đạt $R^2=0.712521$, RMSE $=0.130353$, MAE $=0.085233$ và
MAPE $=7.038\%$. Các chỉ số này không dùng để chọn mô hình chính, nhưng giúp diễn giải sai số theo
đơn vị bán kính Trái Đất.

\begin{figure}[H]
    \centering
    \includegraphics[width=0.82\textwidth]{assets/part2/actual_vs_predicted.png}
    \caption{Actual vs Predicted trên tập test của mô hình Lasso. Các điểm càng gần đường chéo
    thì dự báo càng sát giá trị thật.}
    \label{fig:part2_actual_pred}
\end{figure}

% ------------------------------------------------------------
\subsection{Đánh Giá Mô Hình (Evaluation)}
% ------------------------------------------------------------

\subsubsection{So sánh mô hình tuyến tính}

\begin{table}[H]
\centering
\caption{So sánh các mô hình trên tập test}
\label{tab:part2_model_summary}
\resizebox{\textwidth}{!}{%
\begin{tabular}{@{}lrrrrr@{}}
\hline
\textbf{Mô hình} & \textbf{Test $R^2$} & \textbf{Test RMSE} & \textbf{Test MAE} &
\textbf{CV-MSE} & \textbf{CV-$R^2$} \\
\hline
Lasso & 0.748011 & 0.056917 & 0.037929 & 0.005092 & 0.672730 \\
OLS full & 0.747794 & 0.056941 & 0.037952 & 0.005075 & 0.673407 \\
Ridge & 0.747678 & 0.056954 & 0.038204 & 0.005074 & 0.673433 \\
OLS selected & 0.743913 & 0.057378 & 0.039013 & 0.005103 & 0.671611 \\
\hline
\end{tabular}}
\end{table}

\begin{figure}[H]
    \centering
    \includegraphics[width=\textwidth]{assets/part2/model_comparison.png}
    \caption{So sánh Test $R^2$, RMSE và MAE giữa các mô hình tuyến tính. Lasso đạt Test $R^2$
    cao nhất, nhưng khác biệt giữa Lasso, OLS full và Ridge rất nhỏ.}
    \label{fig:part2_model_comparison}
\end{figure}

Nhận xét:
\begin{itemize}
    \item Lasso có Test $R^2$ cao nhất là 0.748011 và Test RMSE thấp nhất là 0.056917.
    \item OLS full gần như tương đương Lasso, cho thấy quan hệ tuyến tính sau tiền xử lý đã mô tả
          dữ liệu khá tốt.
    \item Ridge có CV-MSE thấp nhất trong nhóm mô hình tuyến tính, phản ánh regularization giúp
          ổn định hơn trên các fold train.
    \item OLS selected ít biến hơn nhưng Test $R^2$ giảm nhẹ xuống 0.743913; việc rút gọn biến giúp
          dễ diễn giải hơn nhưng có thể mất một phần thông tin dự báo.
\end{itemize}

Sự khác biệt giữa Lasso, OLS full và Ridge rất nhỏ: Test $R^2$ của ba mô hình chỉ chênh nhau ở
chữ số thập phân thứ tư. Lý do chính là pipeline đã làm phần việc mà regularization thường hỗ trợ:
log-transform giảm lệch phải, Winsorization giảm ảnh hưởng cực trị, và VIF filtering đã loại
\texttt{log\_pl\_orbper}, \texttt{log\_pl\_eqt} để đưa VIF cuối cùng xuống dưới 10. Khi đa cộng
tuyến đã được xử lý trước, Ridge không còn lợi thế ổn định hệ số quá lớn, còn Lasso với
$\lambda^*=0.0316228$ vẫn không tạo nghiệm thưa. Vì vậy mô hình tuyến tính cơ bản đã gần đạt
hiệu năng tối đa của nhóm mô hình tuyến tính trên tập feature hiện có.

Mặt khác, $R^2_{\mathrm{test}}\approx0.748$ cũng có nghĩa khoảng 25.2\% phương sai của
\texttt{log\_pl\_rade} vẫn chưa được giải thích. Phần còn lại có thể đến từ ba nguồn: quan hệ vật
lý phi tuyến giữa khối lượng--bán kính ở từng cấu tạo hành tinh, thiếu biến mô tả thành phần bên
trong/tuổi hệ sao/độ chính xác phép đo, và selection bias của dữ liệu quan trắc thiên văn. Do đó,
kết quả này nên được đọc như một mô hình nội suy tuyến tính tốt cho nhóm terrestrial/Super-Earth,
không phải một phương trình vật lý đầy đủ của quá trình hình thành hành tinh.

\subsubsection{Phân tích hệ số và feature importance}

\begin{figure}[H]
    \centering
    \includegraphics[width=0.9\textwidth]{assets/part2/feature_importance.png}
    \caption{Feature importance dựa trên hệ số Ridge sau chuẩn hóa. Vì các feature đã được
    standardize, độ lớn hệ số có thể so sánh trực tiếp hơn.}
    \label{fig:part2_feature_importance}
\end{figure}

Các biến quan trọng nhất gồm \texttt{log\_pl\_bmasse}, \texttt{log\_sy\_dist} và
\texttt{log\_pl\_trandep}. Trong đó, \texttt{log\_pl\_bmasse} là biến có ảnh hưởng dương mạnh nhất,
phù hợp với quan hệ vật lý giữa khối lượng và bán kính hành tinh.
\texttt{log\_sy\_dist} cũng có hệ số dương khá lớn, nhưng không nên diễn giải theo nghĩa nhân quả
rằng hành tinh ở xa Trái Đất thì tự nhiên có bán kính lớn hơn. Cách đọc hợp lý hơn là
\emph{selection bias}: ở khoảng cách xa, các kính thiên văn dễ phát hiện và xác nhận những hành
tinh có tín hiệu mạnh hơn, thường là transit sâu hơn hoặc bán kính lớn hơn. Vì vậy hệ số này phản
ánh thiên lệch quan sát của catalog nhiều hơn là một cơ chế vật lý trực tiếp.

\subsubsection{Phân tích phần dư và giả định mô hình}

\begin{figure}[H]
    \centering
    \includegraphics[width=\textwidth]{assets/part2/residual_plots.png}
    \caption{Bốn biểu đồ chẩn đoán phần dư của mô hình Lasso: Residuals vs Fitted, Normal Q-Q,
    Scale-Location và Cook's Distance.}
    \label{fig:part2_residual}
\end{figure}

Nhận xét từ residual plots:
\begin{itemize}
    \item Phần dư tập trung quanh 0, cho thấy mô hình không có thiên lệch hệ thống quá rõ trên đa
          số quan sát.
    \item Q-Q plot cho thấy phần giữa của phân phối phần dư tương đối gần chuẩn, nhưng hai đuôi vẫn
          lệch khỏi đường chéo; đây là dấu hiệu dữ liệu còn ngoại lai hoặc quan hệ phi tuyến.
    \item Cook's Distance lớn nhất được ghi nhận là 0.617979. Với $n=217$ mẫu test, ngưỡng
          $4/n \approx 0.01843$ đánh dấu 17 điểm influential.
\end{itemize}

\begin{table}[H]
\centering
\caption{Kiểm định chẩn đoán trên phần dư test của Lasso}
\label{tab:part2_residual_tests}
\begin{tabular}{@{}lrrp{0.45\textwidth}@{}}
\hline
\textbf{Kiểm định} & \textbf{Statistic} & \textbf{p-value} & \textbf{Diễn giải} \\
\hline
Shapiro-Wilk & 0.887620 & $1.21\times10^{-11}$ &
Bác bỏ giả thuyết phần dư phân phối chuẩn; hai đuôi còn lệch, phù hợp với Q-Q plot. \\
Breusch-Pagan & 36.895113 & 0.000120 &
Bác bỏ giả thuyết phương sai không đổi; vẫn còn heteroskedasticity sau log-transform và
Winsorization. \\
\hline
\end{tabular}
\end{table}

Nhóm không xóa 17 điểm influential vì đây là dữ liệu thiên văn thật, không phải lỗi nhập liệu đã
được chứng minh. Các điểm này có thể đại diện cho hệ hành tinh hiếm, sai số đo lớn, hoặc vùng biên
của catalog. Xóa cứng sẽ làm báo cáo đẹp hơn về mặt diagnostic nhưng dễ tạo thiên lệch mẫu. Vì
vậy nhóm giữ lại chúng, dùng log-transform/Winsorization để giảm ảnh hưởng cực trị, và ghi nhận
đây là một giới hạn khi diễn giải hệ số.

\subsubsection{Kiểm chứng giả thuyết Giãn nở nhiệt (Thermal Expansion Hypothesis)}

\paragraph{Kiểm định tác động tương tác giữa Khối lượng và Nhiệt độ:}
Từ hiện tượng đường LOWESS bị cong nhẹ ở vùng dự đoán cao trên biểu đồ chẩn đoán Residuals vs Fitted, 
giả thuyết được đặt ra là: tác động của khối lượng (lực nén hấp dẫn) lên cấu trúc bán kính hành tinh 
không độc lập, mà bị chi phối bởi mức độ bức xạ nhiệt (hiệu ứng giãn nở). Để kiểm chứng, mô hình hồi 
quy được mở rộng thành một cấu trúc phi tuyến cục bộ thông qua đại lượng tương tác:
\begin{equation}
    Y = \beta_1 X_{\text{mass}} + \beta_2 X_{\text{heat}} + \beta_3 (X_{\text{mass}} \times X_{\text{heat}})
\end{equation}

Thực nghiệm đầu tiên tiến hành bổ sung biến tương tác giữa khối lượng hành tinh và nhiệt độ bề mặt 
của sao chủ (\texttt{log\_pl\_bmasse} $\times$ \texttt{st\_teff}). Kết quả truy xuất cho thấy p-value 
của biến này đạt mức 0.098, vượt ngưỡng mức ý nghĩa $\alpha = 0.05$. Sự thất bại của giả thuyết này 
có thể lý giải bằng bản chất vật lý: nhiệt độ bề mặt sao chủ không phản ánh chính xác thông lượng 
nhiệt (heat flux) mà hành tinh thực nhận nếu khuyết đi tham số về khoảng cách quỹ đạo.

Để khắc phục, thực nghiệm thứ hai thay thế nhiệt độ sao bằng thông lượng bức xạ thực nhận (\texttt{log\_pl\_insol} 
tạo thành biến tương tác (\texttt{log\_pl\_bmasse} $\times$ \texttt{log\_pl\_insol}). Kết quả ghi 
nhận p-value $\approx 0$, khẳng định sự tồn tại của cơ chế tương tác nhiệt - khối lượng trên tập 
huấn luyện. Tuy nhiên, nghịch lý xuất hiện khi kiểm định hiệu năng: hệ số xác định trên tập kiểm thử 
($R^2_{\text{Test}}$) lại sụt giảm từ 0.747 xuống 0.741, đồng thời RMSE tăng.

Bản chất của hiện tượng này là minh chứng kinh điển cho bài toán \emph{Sự đánh đổi giữa Độ chệch và 
Phương sai (Bias-Variance Tradeoff)}. Biến tương tác đã vô tình kích hoạt hiện tượng Đa cộng tuyến 
cấu trúc (Structural Multicollinearity). Mặt phẳng OLS bị ép phải vặn xoắn (giảm Bias) nhằm cố gắng 
nội suy khít với các biến thiên cục bộ sinh ra từ hiệu ứng giãn nở nhiệt trong tập huấn luyện, từ đó 
đánh mất khả năng khái quát hóa quy luật (tăng Variance) khi tiếp xúc với dữ liệu kiểm thử.

\paragraph{Kết luận:}
Nỗ lực triệt tiêu sai số phần dư cục bộ bằng các biến tương tác đã dẫn đến hiện tượng quá khớp (Overfitting). 
Hệ quả này củng cố tính đúng đắn của nguyên lý tinh giản Occam's Razor trong thiết kế Machine Learning: 
cấu trúc OLS thuần túy hiện tại (không chứa đại lượng tương tác) là điểm cân bằng toàn cục tối ưu nhất 
giữa năng lực giải thích hiện tượng và sức mạnh ngoại suy định lượng.

\subsubsection{Phương pháp nâng cao: Kernel Ridge Regression}

Ngoài các mô hình tuyến tính chính, nhóm thử Kernel Ridge Regression với RBF kernel:
\begin{equation}
    k(\mathbf{x},\mathbf{x}')
    = \exp\left(-\gamma\|\mathbf{x}-\mathbf{x}'\|^2\right).
\end{equation}

Do Kernel Ridge cần ma trận Gram kích thước $n\times n$, pipeline dùng tối đa 200 mẫu train để giữ
chi phí tính toán hợp lý. Hyperparameters được chọn trên validation tách từ train, sau đó mới đánh
giá một lần trên test.

\begin{table}[H]
\centering
\caption{Một số cấu hình Kernel Ridge theo validation RMSE}
\label{tab:part2_kernel_search}
\begin{tabular}{@{}rrr@{}}
\hline
\textbf{Alpha} & \textbf{Gamma} & \textbf{Validation RMSE} \\
\hline
0.010 & 0.010 & 0.071636 \\
0.100 & 0.010 & 0.075317 \\
0.001 & 0.010 & 0.075770 \\
0.001 & 0.001 & 0.080698 \\
1.000 & 0.010 & 0.084010 \\
0.010 & 0.001 & 0.089867 \\
0.100 & 0.001 & 0.091168 \\
0.010 & 0.050 & 0.091262 \\
\hline
\end{tabular}
\end{table}

\begin{table}[H]
\centering
\caption{Kết quả Kernel Ridge tốt nhất trên test set}
\label{tab:part2_kernel_result}
\begin{tabular}{@{}lrrrrr@{}}
\hline
\textbf{Mô hình} & \textbf{Alpha} & \textbf{Gamma} & \textbf{Test $R^2$} & \textbf{RMSE} & \textbf{MAE} \\
\hline
Kernel Ridge (RBF) & 0.01 & 0.01 & 0.700766 & 0.062023 & 0.038842 \\
\hline
\end{tabular}
\end{table}

Kết quả Kernel Ridge thấp hơn nhóm mô hình tuyến tính chính trong lần chạy \texttt{part2/real/},
nhưng đây không phải là một so sánh công bằng về năng lực cuối cùng của phương pháp. Giới hạn
200 mẫu train là quyết định kỹ thuật có chủ ý vì Kernel Ridge phải lưu và giải hệ trên ma trận
Gram $n\times n$, làm chi phí bộ nhớ và thời gian tăng xấp xỉ theo $O(n^2)$ đến $O(n^3)$ tùy bước
giải hệ. Do đó, kết quả này chỉ cho thấy cấu hình Kernel Ridge thử nghiệm trong báo cáo chưa vượt
mô hình tuyến tính; nó không đủ để kết luận Kernel Ridge nói chung kém hơn OLS/Ridge/Lasso trên
bài toán này.

\FloatBarrier

% ------------------------------------------------------------
\subsection{Kết Luận Phần 2}
% ------------------------------------------------------------

Phần 2 đã hoàn thành một pipeline data fitting đầy đủ cho dữ liệu thực tế:
\begin{itemize}
    \item Dữ liệu được chọn từ NASA Exoplanet Archive, sau đó giới hạn vào nhóm hành tinh đá và
          Super-Earth với 1084 mẫu.
    \item Báo cáo đã kiểm tra duplicate, missing values, ngoại lai/nhiễu và các khả năng conflict
          trong dữ liệu quan trắc.
    \item Pipeline tiền xử lý gồm log-transform, Winsorization, MICE imputation, VIF filtering và
          standardization; các tham số đều fit trên train để tránh data leakage.
    \item Các mô hình chính tái sử dụng code tự cài ở Phần 1: OLS, Ridge, Lasso, k-fold
          cross-validation và residual analysis.
    \item Trên test set, Lasso là mô hình tuyến tính tốt nhất với Test $R^2=0.748011$,
          RMSE $=0.056917$ và MAE $=0.037929$ trên thang log target.
    \item Khi biến đổi ngược về bán kính gốc, Lasso đạt RMSE $=0.130353$, MAE $=0.085233$ và
          MAPE $=7.038\%$, cho thấy sai số dự báo có thể diễn giải được theo đơn vị bán kính
          Trái Đất.
\end{itemize}

Nhìn chung, báo cáo Part 2 đã bám sát quy trình trong ví dụ tham khảo: có phần giới thiệu bài toán,
hiểu dữ liệu, chuẩn bị dữ liệu, xây dựng mô hình, dự báo và đánh giá mô hình. Đồng thời, so với ví
dụ, phần này bổ sung thêm các kỹ thuật phù hợp với dữ liệu thiên văn như MICE, VIF filtering,
regularization và phân tích feature importance.
